\section{Experiments}
\label{sec:experiments}

\subsection{Experimental Setup}
\label{sec:experimental-setup}

\paragraph{Dataset.} We evaluate on the dataset assembled by
Ibiyo et al.~\cite{ease2025rag}, which mines packages from Datadog's
malicious-software-packages dataset~\cite{datadog-dataset}. The replication
package ships pre-split JSON files: a training partition of 1{,}158 malicious
and 3{,}005 benign packages, and a test partition of 276 malicious and 753
benign packages (1{,}029 in total). Each entry contains the package name,
file list, and the raw \texttt{setup.py} content; test entries additionally
carry an AST-derived textual behavior description. Nine malicious test
entries whose \texttt{setup.py} could not be extracted are excluded from evaluation, leaving 1{,}020 test
samples (267 malicious, 753 benign). The training partition serves as the
knowledge base; the test partition is never used for retrieval. The malicious
class is dominated by typosquatted or squatted packages (e.g.,
\texttt{colorama} clones, \texttt{requests}-lookalikes) whose installation
code frequently hides its malicious behavior in the \texttt{setup.py} entry
point, making the task representative of real-world supply-chain attacks.

\paragraph{Setup.} The seven controlled Llama retrieval methods (Section~\ref{sec:retrieval}) are served
by the same \texttt{Llama-3.1-8B-Instruct} model with identical prompting
and decoding settings, and retrieve the top three documents per query, so
that the only varying factor is the retrieval strategy. Dense indexes use
\texttt{Qwen3-Embedding-4B}; BM25 uses a standard Okapi implementation.
The direct and behavior-guided Jev protocols are specified in
Section~\ref{sec:jev-method}. Neither Jev variant uses retrieval; both are
reported as external model-family comparators rather than as additional Llama
retrieval configurations.

\paragraph{Metrics.} We report precision, recall, and F1 for the malicious
class, benign-class F1, macro-F1, balanced accuracy, false-positive rate
(FPR), and false-negative rate (FNR). Because JSON decoding can fail (e.g.,
output truncation by the serving backend), we additionally report
\emph{coverage}: the fraction of samples whose response parsed successfully.
Classification metrics are conditional on parseable responses, while failures
are reported separately through coverage.
As an approximation of whether the model ``hallucinates'' security-relevant
behaviors, we also ground each claimed behavior in a lightweight static
checker built on Python's \texttt{ast} module, which flags shell execution,
process creation, network access, file writes/deletion, dynamic execution,
environment access, base64 decoding, remote download, and persistence.
Free-form behavior names produced by the model are keyword-normalized to this
vocabulary; a claim is supported when the corresponding static flag is
present on any line of the target. We report the unsupported-claim rate
$= \mathrm{unsupported} / \mathrm{claimed}$, behavior precision, and behavior
recall. Because static analysis is neither sound nor complete on obfuscated
code, we treat these figures as an automated static-grounding proxy,
not as hallucination ground truth. Direct Jev has no generated behavior claims
or citations, so this faithfulness metric is not applicable to it.
Behavior-guided Jev instead yields typed behavior predictions; its comparison
with static flags is reported separately as typed behavior alignment.
For behavior-guided Jev, the same formulas compare its thresholded typed
behavior predictions with the static flags. We call this typed behavior
alignment, rather than explanation faithfulness, because Jev provides neither
free-form behavior explanations nor source citations.
\endinput
